% ============================================================
%  THE ORTYX PROTOCOL
%  Part III — Decomposition
%  Chapter 11: The Audit Operators
% ============================================================

\chapter{The Audit Operators}
\label{ch:audit-operators}

\chapepigraph{A procedure that cannot be stated is not a procedure.
It is an intuition. Intuitions can be correct. But they cannot
be shared, replicated, or applied consistently by people
who did not have them. The transition from intuition to
procedure is the transition from insight to method.}{Flyxion}

\noindent
Parts~I and~II did not produce verdicts. They produced pressure.
The Phoenix corpus was examined at its strongest, followed
through its extensions, and observed accumulating structural
tensions that the framework's own vocabulary was insufficient
to name. The reader who has worked through those ten chapters
should arrive at this point with a specific kind of epistemic
discomfort: the sense that something is structurally wrong
with the framework, combined with the absence of a precise
enough vocabulary to say what it is.

Part~III provides that vocabulary. Not as an external imposition
on the Phoenix corpus but as a natural response to the
demands that Parts~I and~II have generated. The four failure
geometries identified in Part~II --- Decorative Formalism,
Definitional Closure, Semantic Universality Collapse, and
Operational Severance --- are now formalized as components
of a systematic audit procedure. That procedure is named,
for the first time, the \Ortyx{} Protocol.

The naming is deliberate. The protocol has been operating
throughout Parts~I and~II, but it was unnamed. The transition
to Part~III marks the moment at which the protocol recognizes
itself: the moment at which informal diagnostic practice
becomes explicit methodology. This is consistent with the
co-development thesis stated in the Preface. The protocol
does not arrive from outside; it emerges from the demands
of the auditing activity itself.

\section{The System Triple}
\label{sec:system-triple}

The first formal requirement of the \Ortyx{} Protocol is
a notation for the thing being audited. Any theoretical
framework can be described, at the most general level,
as a triple:
\[
  \SysTriple = (\mathcal{A}, \mathcal{M}, \mathcal{E}),
\]
where $\mathcal{A}$ is the set of axioms or foundational
commitments, $\mathcal{M}$ is the set of mechanisms or
processes the framework proposes, and $\mathcal{E}$ is
the set of empirical claims the framework makes about
the world.

This representation is deliberately coarse. It does not
capture the full structure of a theoretical framework;
no compact formal representation can. Its purpose is
to make three structural features of any framework
explicit: what it takes for granted ($\mathcal{A}$),
what it proposes happens ($\mathcal{M}$), and what it
claims is true of the observable world ($\mathcal{E}$).
The distinction between these three components is not
always maintained in the literature, and one of the
recurring sources of the failure geometries in Part~II
is the migration of claims between components without
acknowledgment.

For the Phoenix corpus, a preliminary triple can be stated:

$\mathcal{A}$: Temporal coherence is the primary organizing
principle of meaningful information; biological sensory
systems have evolved to prioritize temporally coherent
signals; the auditory system is an appropriate model
for general cognitive processing.

$\mathcal{M}$: The Marine salience filter detects temporal
coherence; the \MEMeight{} field encodes salient events
as wave functions; the \HCFone{} conformity relation
compresses harmonic structure relationally; the three
systems compose into a salience-centered cognitive
architecture.

$\mathcal{E}$: Lossy compression increases Marine-detectable
jitter at block boundaries; harmonic consensus encoding
achieves lower bit rates than independent encoding on
strongly coupled sources; the \MEMeight{} architecture
produces more efficient associative retrieval than
static vector stores; systems trained on salience-filtered
representations outperform dense systems on specified
benchmarks.

\begin{technote}[Technical Note: The Audit Triple]
The system triple $\SysTriple = (\mathcal{A}, \mathcal{M},
\mathcal{E})$ makes one structural requirement explicit:
$\mathcal{E}$ must be non-empty and its elements must
be independently accessible --- meaning, the truth value
of empirical claims must be determinable by procedures
that do not presuppose $\mathcal{A}$. A framework in
which all of $\mathcal{E}$ is entailed by $\mathcal{A}$
alone, without contribution from $\mathcal{M}$, exhibits
$\FailGeo{2}$: its empirical claims are definitional
consequences of its axioms rather than independent
discoveries about the world.
\end{technote}

\section{The Four Audit Operators}
\label{sec:four-operators}

The \Ortyx{} Protocol applies four operators to a system
triple $\SysTriple$. Each operator returns a score in
$[0,1]$, where higher scores indicate greater audit
compliance. The global epistemic stability functional
$\EpStab(\SysTriple)$ is then a weighted combination
of the four scores:
\[
  \EpStab(\SysTriple) = \alpha\,\AuditS(\SysTriple)
  + \beta\,\AuditF(\SysTriple)
  - \gamma\,\AuditP(\SysTriple)
  - \delta\,\AuditR(\SysTriple),
\]
where the positive terms represent structural virtues
and the negative terms represent structural liabilities.
The weights $\alpha, \beta, \gamma, \delta > 0$ are
context-sensitive parameters that reflect the relative
importance of each dimension for the framework under
evaluation. Their determination is discussed in Section
\ref{sec:weight-determination}.

The four operators are:

\subsection{The Structural Consistency Operator \texorpdfstring{$\AuditS$}{S}}
\label{subsec:structural-consistency}

$\AuditS(\SysTriple)$ measures the degree to which the
claims in $\mathcal{E}$ are non-trivially derived from
the combination of $\mathcal{A}$ and $\mathcal{M}$,
rather than either (i)~entailed by $\mathcal{A}$ alone
(indicating $\FailGeo{2}$) or (ii)~independent of
$\mathcal{A}$ and $\mathcal{M}$ (indicating disconnection
between the framework's theoretical machinery and its
empirical claims).

A high $\AuditS$ score indicates that the framework's
mechanisms do real explanatory work: that the empirical
claims follow from the mechanisms applied to the axioms
in non-trivial ways that generate predictions beyond
what the axioms alone would support.

Formally: let $\mathcal{E}_{\mathcal{A}}$ denote the
claims in $\mathcal{E}$ entailed by $\mathcal{A}$ alone,
and $\mathcal{E}_{\mathcal{M}}$ denote the claims derivable
from $\mathcal{M}$ combined with $\mathcal{A}$. Then:
\[
  \AuditS(\SysTriple)
  = \frac{|\mathcal{E}_{\mathcal{M}} \setminus
    \mathcal{E}_{\mathcal{A}}|}{|\mathcal{E}|},
\]
the fraction of empirical claims that are mechanistic
consequences beyond what the axioms alone predict.
High values indicate that the mechanisms are doing
real inferential work.

\subsection{The Falsifiability Operator \texorpdfstring{$\AuditF$}{F}}
\label{subsec:falsifiability}

$\AuditF(\SysTriple)$ measures the degree to which
the framework specifies, for each claim in $\mathcal{E}$,
an observation that would count as disconfirmation.
This is not the narrow Popperian requirement that a
single claim must be falsifiable in isolation; it is
the more practical requirement that the framework's
practitioners can specify, for each empirical claim,
what range of observations they would interpret as
evidence against it.

A low $\AuditF$ score can arise from any of the
failure geometries: $\FailGeo{1}$ (claims expressed
in notation too underspecified to generate observations);
$\FailGeo{2}$ (claims definitionally guaranteed true);
$\FailGeo{3}$ (claims broadened to accommodate all
possible observations); $\FailGeo{4}$ (claims expressed
in terms of quantities with no measurement procedure).

Operationally, $\AuditF$ is assessed by the following
procedure: for each claim $e \in \mathcal{E}$, ask
the question ``What observation $O$ would constitute
evidence against $e$?'' The response admits three
classifications: (i)~a specific, observable $O$ is
named (the claim is falsifiable); (ii)~no specific
$O$ is named but the framework's practitioners agree
one could in principle be specified (the claim is
potentially falsifiable); (iii)~any candidate $O$
can be accommodated by the framework (the claim is
unfalsifiable within the framework). $\AuditF$ is
the weighted proportion of claims in category~(i)
and category~(ii).

\subsection{The Projection Dependence Operator \texorpdfstring{$\AuditP$}{P}}
\label{subsec:projection-dependence}

$\AuditP(\SysTriple)$ measures the degree to which
the framework's claims depend on projection operators
that lose information in ways that cannot be detected
from within the projected representation. High
$\AuditP$ scores indicate that the framework is making
claims about a reduced representation $M$ as if they
were claims about the underlying space $X$, without
establishing that the projection $\pi: X \to M$
is constraint-faithful.

Formally, the projection defect of a framework is:
\[
  \ProjDef(\SysTriple) = \sup_{x \in X_{\mathrm{adm}}}
  \abs{\ConstrX(x) - \ConstrM(\pi(x))},
\]
where $\ConstrX$ is the constraint operator on the
full space and $\ConstrM$ is the induced constraint
on the reduced representation. A framework with high
$\AuditP$ score is one whose projection defect is
large: its claims about the representation do not
reliably track the constraints operative in the
underlying space.

For the Phoenix corpus, the most significant projection
dependence arises in the move from waveform properties
(the full space) to Marine jitter metrics (the
reduced representation). The claim that jitter
detects ``meaningful structure'' is a claim about
the full space (the signal has meaningful structure)
made on the basis of a projection (the jitter
metric). Whether this projection is constraint-faithful
--- whether signals with low jitter genuinely have
meaningful structure in any sense independent of
the framework's definition --- is exactly what
$\AuditP$ measures.

\subsection{The Recursive Closure Operator \texorpdfstring{$\AuditR$}{R}}
\label{subsec:recursive-closure}

$\AuditR(\SysTriple)$ measures the degree to which
the framework has become self-referentially closed:
the degree to which its own vocabulary is the
primary resource for evaluating its claims, with
insufficient purchase available for external
correction. Formally, this is related to the
semantic closure metric:
\[
  \SemClose(t)
  = \frac{\norm{F(M_t)}}{\norm{G(X_t)}},
\]
where $F(M_t)$ represents the evolution of the
representation under internal dynamics and $G(X_t)$
represents the correction available from the underlying
system. As $\SemClose \to \infty$, the representation
evolves primarily under its own dynamics, with
decreasing purchase from external reality.

High $\AuditR$ scores indicate that the framework
has developed internal mechanisms for accommodating
potential disconfirmations --- through vocabulary
migration, definitional expansion, or reinterpretation
strategies --- in ways that insulate it from the
corrective pressure that would otherwise force revision.

\section{Applying the Operators to the Phoenix Corpus}
\label{sec:applying-operators}

With the operators defined, they can be applied
systematically to the Phoenix triple.

\subsection{Structural Consistency}

The Phoenix corpus's $\AuditS$ score is highest for
the engineering components. The claim that lossy
compression increases Marine-detectable jitter is
a non-trivial consequence of the Marine mechanism
applied to the axiom that block-based transform
coding introduces block boundary discontinuities.
The claim that harmonic consensus encoding reduces
bit rate is a non-trivial consequence of the \HCFone{}
mechanism applied to the axiom that strongly coupled
harmonic sources have high conformity rates.

The $\AuditS$ score drops for the architectural and
philosophical claims. The claim that temporal coherence
constitutes the organizing principle of consciousness
is not a non-trivial consequence of the Marine and
\MEMeight{} mechanisms; it is a claim at a different
level that the mechanisms cannot reach without
additional bridging principles that the corpus has
not established.

\begin{auditprop}
\label{ap:structural-consistency-result}
The Phoenix corpus exhibits high structural consistency
($\AuditS$ close to 1.0) for its engineering-level
claims ($\mathcal{E}_{\mathrm{eng}}$) and substantially
lower structural consistency ($\AuditS < 0.5$) for
its architectural and philosophical claims
($\mathcal{E}_{\mathrm{arch}}, \mathcal{E}_{\mathrm{phil}}$).
The mechanisms in $\mathcal{M}$ do real inferential
work at the engineering level but do not reach the
philosophical claims without steps that are either
axiomatic imports or definitional moves.
\end{auditprop}

\subsection{Falsifiability}

The Phoenix corpus's $\AuditF$ score is highest for
the jitter and conformity claims. The prediction that
Marine-derived jitter is higher in compressed audio
than in the original at block boundaries is specific
enough that an experiment could confirm or deny it.
The conformity rate prediction is similarly specific.

The $\AuditF$ score drops sharply for the consciousness
and authenticity claims. As Chapter~9 established,
the temporal coherence concept has been broadened to
the point where both coherent and incoherent salience
can be accommodated. The authenticity concept has been
noted as potentially functioning as a definitional
label rather than an independently measurable property.
For these claims, no specific disconfirming observation
has been specified.

\subsection{Projection Dependence}

The Phoenix corpus's $\AuditP$ score is elevated
throughout. The fundamental move of the corpus ---
inferring meaningful structure from signal properties
--- is a projection-dependent inference that the corpus
has not established as constraint-faithful. The most
significant projection dependence occurs in the
consciousness claims, where properties of the \MEMeight{}
field (a formal representation) are attributed to
conscious experience (the underlying space) without
establishing the constraint relationship between them.

\subsection{Recursive Closure}

The Phoenix corpus's $\AuditR$ score is moderate and
growing. The vocabulary migration of ``coherence'' across
levels (Chapter~6) is one mechanism of recursive closure:
observations in any domain can be redescribed using
the coherence vocabulary in ways that make them appear
confirmatory. The accommodation strategy for unexpected
salience (Chapter~9) is another: any observation can
be reinterpreted as a coherence phenomenon. These are
the beginning of closure dynamics, not yet runaway
closure, which is why the score is moderate rather
than maximal.

\section{The Weight Determination Problem}
\label{sec:weight-determination}

The global epistemic stability functional $\EpStab$ depends
on weights $\alpha, \beta, \gamma, \delta$ that are not
determined by the formal specification of the operators.
This is a genuine limitation of the framework, and it
deserves explicit acknowledgment.

The weights reflect context-sensitive judgments about
which epistemic virtues are most important for a given
framework at a given stage of development. For a framework
at an early stage, where mechanisms are being proposed
and evidence is not yet available, a high weight on
$\AuditF$ (falsifiability) relative to the others is
appropriate: the framework should be generating testable
predictions even if they have not yet been tested.
For a framework at a mature stage, where extensive
empirical work has been done, a high weight on
$\AuditS$ (structural consistency) reflects the
importance of mechanisms actually deriving the
observed results rather than merely accommodating them.

This weight-sensitivity is not a bug in the audit
protocol; it is a feature. Different stages of
scientific development warrant different epistemic
standards, and an audit protocol that applied
identical standards across all stages would either
be too demanding for early-stage proposals or too
permissive for mature theories. The Ortyx Protocol
makes the weight-sensitivity explicit rather than
hiding it.

\begin{auditprop}
\label{ap:context-sensitivity}
The \Ortyx{} audit is context-sensitive with respect
to the weights $\alpha, \beta, \gamma, \delta$ of
the epistemic stability functional. For a given
framework, the appropriate weights depend on: (i)~the
stage of development of the framework, (ii)~the
domain in which it operates, and (iii)~the purposes
for which the audit is conducted. An audit protocol
that does not acknowledge its own weight-sensitivity
is itself exhibiting a form of Definitional Closure:
it has built its evaluative conclusions into its
parameter choices without acknowledging that those
choices could be otherwise.
\end{auditprop}

\section{The Decomposition}
\label{sec:decomposition}

With the operators applied, the formal decomposition
of the Phoenix corpus can be stated:

\audittrace{\SysTriple_{\mathrm{Phoenix}}}{\ThetaS}{\ThetaU}

The survivable component $\ThetaS$ contains:

\begin{itemize}
  \item The Marine jitter metric and its engineering
  applications: fully operationalized, with testable
  predictions about jitter distributions in compressed
  versus uncompressed audio.
  \item The harmonic conformity relation and the
  consensus encoding architecture: engineering-level
  claims with a sound theoretical basis and testable
  efficiency predictions.
  \item The associative memory efficiency of wave
  superposition: established in the holographic memory
  literature, with the \MEMeight{} formalism providing
  a specific parameterization that admits comparison
  with alternatives.
  \item The predictive processing alignment: the
  structural parallel between salience-gated encoding
  and predictive processing error signals, which
  motivates a research connection to the neuroscience
  literature.
  \item The behavioral injection / context poisoning
  framework: at \LevelI{}, as a description of a
  real vulnerability class in current AI systems.
\end{itemize}

The unstable component $\ThetaU$ contains:

\begin{itemize}
  \item The consciousness claims: not derivable from
  the engineering mechanisms, operationally severed,
  and exhibiting high projection dependence.
  \item The authenticity claims: exhibiting Definitional
  Closure risk, operationally severed, and broadened
  to near-universal accommodation.
  \item The strong AI architecture claims: the assertion
  that the salience-centered paradigm constitutes a
  general foundation for machine intelligence, rather
  than a promising research programme in specific domains.
  \item The Ultrasonic Consciousness Hypothesis:
  empirically contested and operationally severed
  in its Phoenix formulation.
\end{itemize}

\medskip

This decomposition is not a final verdict. It is a
structural diagnosis at the current state of the
framework's development. Some of what is currently in
$\ThetaU$ may migrate to $\ThetaS$ if the framework
provides the operationalization and falsification
structure that currently separates them. The Interlude
between Parts~III and~IV will examine what form that
migration would take.

Chapter~12 applies the decomposition to the
individual failure geometries systematically, providing
a case-by-case audit of the Phoenix corpus against
the formal criteria developed in this chapter.
